Simétricamente, una estructura de extensión en la dirección de $Q$,
es decir, \sisi{de $P_Q$ por $G_Q$}{de $Q_P$ por $G_P$}, se describe
mediante isomorfismos de torsores bajo $G_S$
\begin{equation}\tag{2.0.7}
 \psi_{p;q,q'}:E_{p,q}E_{p,q'}
 \xrightarrow{\ \sim\ }\sisi{E_{pp',q}}{E_{p,qq'}},
\end{equation}
para $p\in P(S)$ y $q,q'\in Q(S)$. Estos son funtoriales en $S$ y
satisfacen las condiciones simétricas de asociatividad y
conmutatividad
\begin{equation}\tag{2.0.8}
\begin{tikzcd}[column sep=5.1em,row sep=3.4em]
 & E_{p,qq'}E_{p,q''}
   \arrow[rd,"{\psi_{p;qq',q''}}"] & \\
 E_{p,q}E_{p,q'}E_{p,q''}
   \arrow[ru,"{\psi_{p;q,q'}\wedge\mathrm{id}}"]
   \arrow[rd,"{\mathrm{id}\wedge\psi_{p;q',q''}}"']
 && E_{p,qq'q''} \\
 & E_{p,q}E_{p,q'q''}
   \arrow[ru,"{\sisi{\psi_{p;q;q'q''}}{\psi_{p;q,q'q''}}}"'] &
\end{tikzcd}
\end{equation}
para $p\in P(S)$ y $q,q',q''\in Q(S)$, y
\begin{equation}\tag{2.0.9}
\begin{tikzcd}[column sep=5em,row sep=3.1em]
 E_{p,q}E_{p,q'} \arrow[r,"{\psi_{p;q,q'}}"]
   \arrow[d,"\mathrm{sym}"']
 & E_{p,qq'} \arrow[d,"\mathrm{id}"] \\
 E_{p,q'}E_{p,q} \arrow[r,"{\psi_{p;q',q}}"']
 & E_{p,q'q}.
\end{tikzcd}
\end{equation}
Aquí $p\in P(S)$ y $q,q'\in Q(S)$.

%%%%%%%%%%  scan idx 170  |  folio 159  %%%%%%%%%%

\underline{Definición 2.1. Sean $P,Q,G$ grupos abelianos de
$\underline T$, y sea $E$ un torsor bajo $G_{P\times Q}$ provisto de
una estructura de extensión en la dirección de $P$, es decir,
\sisi{de $Q_P$ por $G_P$}{de $P_Q$ por $G_Q$}, definida por los
isomorfismos (2.0.4), y de una estructura de extensión en la dirección
de $Q$, es decir, \sisi{de $P_Q$ por $G_Q$}{de $Q_P$ por $G_P$},
definida por los isomorfismos (2.0.7), ambas funtoriales en $S$ y
sujetas a las condiciones de asociatividad y conmutatividad anteriores.
Se dice que las dos estructuras son compatibles si, para todo objeto
$S$ de $\underline T$ y $p,p'\in P(S)$, $q,q'\in Q(S)$, conmuta el
diagrama siguiente:}
\begin{equation}\tag{2.1.1}
\begin{tikzcd}[column sep=2.7em,row sep=3.2em,font=\scriptsize]
 && E_{p,qq'}E_{p',qq'}
    \arrow[rrdd,"{\varphi_{p,p';qq'}}"] && \\
 E_{p,q}E_{p,q'}E_{p',q}E_{p',q'}
    \arrow[rru,"{\psi_{p;q,q'}\wedge\psi_{p';q,q'}}"]
    \arrow[d,"{}"'] &&&& \\
 E_{p,q}E_{p',q}E_{p,q'}E_{p',q'}
    \arrow[rrd,"{\varphi_{p,p';q}\wedge\varphi_{p,p';q'}}"']
 &&&& E_{pp',qq'} \\
 && E_{pp',q}E_{pp',q'}
    \arrow[rru,"{\psi_{pp';q,q'}}"'] &&
\end{tikzcd}
\end{equation}
\underline{donde la flecha vertical se deduce del isomorfismo de
simetría}
\[
 E_{p,q'}E_{p',q}\xrightarrow{\ \sim\ }E_{p',q}E_{p,q'}.
\]
\underline{Se llama \emph{biextensión} de $(P,Q)$ por $G$ a un torsor
$E$ bajo $G_{P\times Q}$ provisto de una estructura de extensión
\sisi{de $Q_P$ por $G_P$}{de $P_Q$ por $G_Q$} y de una estructura de
extensión \sisi{de $G_Q$ por $P_Q$}{de $Q_P$ por $G_P$}, compatibles
en el sentido precedente.}

%%%%%%%%%%  scan idx 171  |  folio 160  %%%%%%%%%%

2.1.1. Equivalentemente, una estructura de biextensión sobre el torsor
$E$ bajo $G_{P\times Q}$ viene dada por dos leyes de composición
parcialmente definidas $\varphi(x,y)$ y $\psi(x,y)$ sobre $E$. La
primera está definida cuando $\pr_2j(x)=\pr_2j(y)$, y la segunda cuando
$\pr_1j(x)=\pr_1j(y)$, donde $j:E\longrightarrow P\times Q$ es el
morfismo estructural del torsor. Además de relaciones de asociatividad
del tipo
\begin{equation}\tag{2.1.1.1}
 \varphi(gx,y)=\varphi(x,gy)=g\varphi(x,y)
 \qquad(x,y\in E(S),\ g\in G(S))
\end{equation}
y de las relaciones análogas para $\psi$, los axiomas dicen que estas
leyes de composición parcial son asociativas y conmutativas y que
satisfacen
\begin{equation}\tag{2.1.1.2}
 \varphi(\psi(x,y),\psi(z,t))
 =\psi(\varphi(x,z),\varphi(y,t)),
\end{equation}
donde $x,y,z,t$ pertenecen respectivamente a
$E_{p,q}(S),E_{p,q'}(S),E_{p',q}(S),E_{p',q'}(S)$.

